% !TEX root = ../algebra_02.tex

\section{Единственность поля из $p^n$ элементов}

\begin{Theorem}{Единственность конечного поля с точностью до изоморфизма}{}
	Пусть $K$ и $L$ --- конечные поля из $p^n$ элементов. Тогда $K \cong L$.
\end{Theorem}

\begin{proof}
	Поле $K$ можно представить в виде $K = \mathbb{F}_p[x]/(f)$, где $f \in \mathbb{F}_p[x]$ --- неприводимый многочлен степени $n$ (он является минимальным многочленом для примитивного элемента поля $K$).

	Заметим, что $f \mid (X^{p^n} - X)$. Действительно, порядок мультипликативной группы равен $|K^*| = p^n - 1$. Для любого $a \in K^*$ выполнено $a^{p^n - 1} = 1$. Умножая на $a$, получаем, что для любого $a \in K$ выполнено $a^{p^n} - a = 0$. В частности, это верно для $x = \bar{X} \in K$, то есть многочлен $X^{p^n} - X$ делится на минимальный многочлен $f$ элемента $x$.

	В кольце многочленов $L[x]$ многочлен $X^{p^n} - X$ имеет ровно $p^n$ корней (поскольку для любого элемента $a \in L$ выполнено $a^{p^n} = a$). Следовательно, $X^{p^n} - X$ полностью раскладывается в $L[x]$ на линейные множители. Из делимости следует, что $f$ также раскладывается на линейные множители над $L$.

	Пусть $b \in L$ --- некоторый корень многочлена $f$. Рассмотрим простое расширение $\mathbb{F}_p(b) / \mathbb{F}_p$. Минимальным многочленом элемента $b$ над $\mathbb{F}_p$ является $f$, поэтому $[\mathbb{F}_p(b) : \mathbb{F}_p] = \deg f = n$. Следовательно, $|\mathbb{F}_p(b)| = p^n$.

	Так как $\mathbb{F}_p(b) \subset L$ и оба поля состоят из $p^n$ элементов, то $\mathbb{F}_p(b) = L$.
	Итоговый изоморфизм строится следующим образом:
	\[ L = \mathbb{F}_p(b) \cong \mathbb{F}_p[x]/(f) = K \]
\end{proof}

\begin{Remark}{Обозначение конечного поля}{}
	Поскольку конечные поля одного порядка изоморфны, поле из $p^n$ элементов обозначается стандартно как $\mathbb{F}_{p^n}$.
\end{Remark}

\begin{Example}{Изоморфизм полей из 8 элементов}{}
	Рассмотрим два представления поля из 8 элементов: $\mathbb{F}_2[x]/(x^3+x+1)$ и $\mathbb{F}_2[y]/(y^3+y^2+1)$. Изоморфизм между ними задается отображением $x \mapsto y+1$, поскольку $y+1$ является корнем многочлена $X^3+X+1$ во втором поле.
\end{Example}
